% =========================== 
% Core encoding & layout 
% =========================== 
\usepackage{geometry} 
\usepackage[indent]{parskip} 
\geometry{left=1in,right=1in,top=1in,bottom=1in} 
\setlength{\parskip}{1em} 

% =========================== 
% Fonts & micro-typography (LuaLaTeX) 
% =========================== 
\usepackage{fontspec} 
\defaultfontfeatures{Ligatures=TeX} 
% \usepackage{pagella-otf}
% \setmainfont{pagella-otf}
\usepackage[nomath]{libertinus}
\usepackage{menukeys}

% Palatino-like text font for LuaLaTeX (TeX Gyre Pagella is the standard free Palatino clone) 
% \setmainfont{TeX Gyre Pagella} 

\usepackage{microtype} 

% =========================== 
% Graphics & color 
% =========================== 
\usepackage{graphicx} 
\usepackage[table]{xcolor} 
\graphicspath{{figures/}{../figures/}} 
\definecolor{Green}{HTML}{1B7F5A} 
\definecolor{Defns}{HTML}{CC7E7E} %Pink 
\definecolor{ERs}{HTML}{7DA2C7} %Blue 
\definecolor{Ps}{HTML}{A17800} %Gold 

% =========================== 
% Math & diagrams 
% =========================== 
\usepackage{amsmath,amsthm} 
\usepackage{amscd} 
\usepackage{amssymb}
\usepackage{mathtools} 
% If you want matching math font in a Palatino style, use unicode-math: 
% (Recommended, but note it changes the math font stack vs pdfLaTeX.) 
\usepackage{unicode-math}
\setmathfont{STIX Two Math}
\setmathfont[range=\mathcal]{STIX Two Math}

% \usepackage{mathabx}
\usepackage{tikz} 
\usepackage{tikz-cd} 
\usepackage{bbm} 
\usepackage{extarrows}

% ===========================
%  Preferred symbol variants / defaults
% ===========================
\let\emptyset\varnothing
\let\leq\leqslant
\let\le\leqslant
\let\geq\geqslant
\let\ge\geqslant

\renewcommand{\iff}{\;\Longleftrightarrow\;}

% ===========================
%  Hyperlinks (color + underline)
% ===========================
\usepackage{xurl}
\usepackage[hidelinks]{hyperref}
\hypersetup{
  colorlinks=true,
  linkcolor=Green,
  citecolor=Green,
  urlcolor=Green,
  pdfborderstyle={/S/U/W 1}
}
% Optional hard underline fallback:
% \usepackage[normalem]{ulem}
% \let\OldHref\href
% \renewcommand{\href}[2]{\OldHref{#1}{\uline{#2}}}

% ===========================
%  Numbered items: counters
%  - Main stream: thmsec (shared by everything except exercises)
%  - Exercise stream: exsec (independent)
%  Format: <section>.<count>
% ===========================
\newcounter{thmsec}[section]
\renewcommand{\thethmsec}{\thesection.\arabic{thmsec}}

\newcounter{exsec}[section]
\renewcommand{\theexsec}{\thesection.\arabic{exsec}}

% ===========================
%  Boxed theorem environments (B) via tcolorbox
% ===========================
\usepackage[most]{tcolorbox}
\tcbuselibrary{breakable,skins}

% --- Colors ---
\definecolor{DefBody}{HTML}{FFFACD}
\definecolor{DefTitle}{HTML}{FFF6A8}
\definecolor{ThmBody}{HTML}{E2E8F4}
\definecolor{ThmTitle}{HTML}{B7C8E6}

% --- Base style for all theorem boxes ---
\tcbset{
  mythm/.style={
    enhanced,
    breakable,
    boxrule=0.9pt,
    arc=5mm,
    coltitle=black,
    fonttitle=\bfseries,
    left=4mm,right=4mm,top=3mm,bottom=3mm,
    boxed title style={boxrule=0pt,arc=5mm},
    titlerule=0.9pt,
    toptitle=2mm,bottomtitle=2mm,
    lefttitle=4mm,righttitle=4mm,
  }
}

% ===========================
%  Internal helper: start numbered item (boxed/unboxed)
%  #1 = label, #2 = printed type (e.g. Definition), #3 = title text
%  - No reservation system: increments the appropriate counter at print time
%  - Creates hypertarget + \label{<label>} at print time
% ===========================
\makeatletter
\newcommand{\BeginNumberedItem}[3]{%
  \refstepcounter{thmsec}%
  \edef\ItemNumber{\thethmsec}%
  \hypertarget{#1}{}\label{#1}%
  \def\ItemType{#2}%
  \def\ItemTitle{#3}%
}
% Same interface, but uses the exercise counter
\newcommand{\BeginExerciseItem}[3]{%
  \refstepcounter{exsec}%
  \edef\ItemNumber{\theexsec}%
  \hypertarget{#1}{}\label{#1}%
  \def\ItemType{#2}%
  \def\ItemTitle{#3}%
}
\makeatother

% ===========================
%  Box templates (no counters inside tcolorbox)
%  IMPORTANT: now uses \ItemNumber (set by begin-helper)
% ===========================
\newtcolorbox{DefBox}[0]{%
  mythm,
  colback=DefBody,
  colbacktitle=DefTitle,
  colframe=black,
  title={\ItemType~\ItemNumber:~\ItemTitle},
}

\newtcolorbox{ThmBox}[0]{%
  mythm,
  colback=ThmBody,
  colbacktitle=ThmTitle,
  colframe=black,
  title={\ItemType~\ItemNumber:~\ItemTitle},
}

\newtcolorbox{ExBox}[0]{%
  mythm,
  colback=white,
  colbacktitle=black!10,
  colframe=black,
  title={\ItemType~\ItemNumber:~\ItemTitle},
}

\newtcolorbox{NoteBox}[0]{%
  mythm,
  colback=black!2,
  colbacktitle=black!12,
  colframe=black!35,
  title={\ItemType~\ItemNumber:~\ItemTitle},
}

% ===========================
%  Public boxed environments (B): {title}{label}
% ===========================
\newenvironment{definitionB}[2]{%
  \BeginNumberedItem{#2}{Definition}{#1}%
  \begin{DefBox}%
}{%
  \end{DefBox}%
}

\newenvironment{exerciseB}[2]{%
  \BeginExerciseItem{#2}{Exercise}{#1}%
  \begin{ExBox}%
}{%
  \end{ExBox}%
}

\newenvironment{theoremB}[2]{%
  \BeginNumberedItem{#2}{Theorem}{#1}%
  \begin{ThmBox}%
}{%
  \end{ThmBox}%
}

\newenvironment{propositionB}[2]{%
  \BeginNumberedItem{#2}{Proposition}{#1}%
  \begin{ThmBox}%
}{%
  \end{ThmBox}%
}

\newenvironment{corollaryB}[2]{%
  \BeginNumberedItem{#2}{Corollary}{#1}%
  \begin{ThmBox}%
}{%
  \end{ThmBox}%
}

\newenvironment{conjectureB}[2]{%
  \BeginNumberedItem{#2}{Conjecture}{#1}%
  \begin{ThmBox}%
}{%
  \end{ThmBox}%
}

\newenvironment{axiomB}[2]{%
  \BeginNumberedItem{#2}{Axiom}{#1}%
  \begin{ThmBox}%
}{%
  \end{ThmBox}%
}

\newenvironment{observationB}[2]{%
  \BeginNumberedItem{#2}{Observation}{#1}%
  \begin{NoteBox}%
}{%
  \end{NoteBox}%
}

% ===========================
%  Deferred/unboxed environments (R): {label}
%  - Same numbering <section>.<count>
%  - Exercises use exsec; everything else uses thmsec
%  - Proofs: unnumbered
% ===========================
\newcommand{\RHeading}[2]{%
  \par\medskip\noindent\textbf{#1~\ItemNumber.}\ %
  \if\relax\detokenize{#2}\relax\else\textbf{#2}\ \fi
}

\newenvironment{exerciseR}[1]{%
  \BeginExerciseItem{#1}{Exercise}{}%
  \RHeading{Exercise}{}%
}{\par\medskip}

\newenvironment{lemmaR}[1]{%
  \BeginNumberedItem{#1}{Lemma}{}%
  \RHeading{Lemma}{}%
}{\par\medskip}

\newenvironment{propositionR}[1]{%
  \BeginNumberedItem{#1}{Proposition}{}%
  \RHeading{Proposition}{}%
}{\par\medskip}

\newenvironment{corollaryR}[1]{%
  \BeginNumberedItem{#1}{Corollary}{}%
  \RHeading{Corollary}{}%
}{\par\medskip}

\newenvironment{axiomR}[1]{%
  \BeginNumberedItem{#1}{Axiom}{}%
  \RHeading{Axiom}{}%
}{\par\medskip}

\newenvironment{assumptionR}[1]{%
  \BeginNumberedItem{#1}{Assumption}{}%
  \RHeading{Assumption}{}%
}{\par\medskip}

\newenvironment{exampleR}[1]{%
  \BeginNumberedItem{#1}{Example}{}%
  \RHeading{Example}{}%
}{\par\medskip}

\newenvironment{remarkR}[1]{%
  \BeginNumberedItem{#1}{Remark}{}%
  \RHeading{Remark}{}%
}{\par\medskip}

% Proofs: unboxed, unnumbered, but still labeled/anchored if you want
\newenvironment{proofR}[1]{%
  \hypertarget{#1}{}\label{#1}%
  \begin{proof}%
}{%
  \end{proof}%
}

% Convenience (for already-defined labels)
\newcommand{\laterref}[2]{\hyperlink{#1}{#2~\ref{#1}}}

% ===========================
%  Mathematical operators
% ===========================

% Linear algebra
\DeclareMathOperator{\spn}{span}
\DeclareMathOperator{\range}{range}
\DeclareMathOperator{\im}{im}
\DeclareMathOperator{\nullspace}{null}
\DeclareMathOperator{\hess}{Hess}

% Multilinear algebra
\DeclareMathOperator{\sym}{sym}
\DeclareMathOperator{\alt}{alt}
\DeclareMathOperator{\sgn}{sgn}
\DeclareMathOperator{\vol}{vol}
\DeclareMathOperator{\tr}{tr}

% Set theory / Topology

\DeclareMathOperator{\inte}{int}
\DeclareMathOperator{\ext}{ext}

% ===========================
%  Commands
% ===========================

% Mathbb
\newcommand{\RR}{\mathbb R}
\newcommand{\CC}{\mathbb C}
\newcommand{\FF}{\mathbb F}

% Mathcal
\newcommand{\LL}{\mathcal L}

% Formatting and notation
\newcommand{\bb}[1]{\symbf{#1}}
\newcommand{\bbi}[1]{\symfit{#1}}
\renewcommand{\rm}[1]{\symrm{#1}}
\renewcommand{\v}[1]{\symbfit{#1}}
\DeclareMathOperator{\supp}{supp}

% Delimiters
\DeclarePairedDelimiter{\paren}{(}{)}
\DeclarePairedDelimiter{\inner}{\langle}{\rangle}
\DeclarePairedDelimiter{\bars}{|}{|}
\DeclarePairedDelimiter{\set}{\{}{\}}

% -- Defining norm with subscript --
\DeclarePairedDelimiter{\normdelim}{\lVert}{\rVert}

% syntax: \norm{X}[p] or \norm*{X}[p] or \norm[\Big]{X}[p]
\NewDocumentCommand{\norm}{s o m o}{%
  \IfBooleanTF{#1}{%
    \IfNoValueTF{#4}
      {\normdelim*{#3}}
      {{\normdelim*{#3}}_{#4}}%
  }{%
    \IfNoValueTF{#2}{%
      \IfNoValueTF{#4}
        {\normdelim{#3}}
        {{\normdelim{#3}}_{#4}}%
    }{%
      \IfNoValueTF{#4}
        {\normdelim[#2]{#3}}
        {{\normdelim[#2]{#3}}_{#4}}%
    }%
  }%
}